%% ============================================================================
%% Fichier  : annexes/Ann_C_templates.tex
%% Titre    : Modèles et Grilles
%% Version  : 1.0 -- novembre 2026 (INF3226, d'après ACSI)
%% ============================================================================
%% Outils O0 à O9 (Chapitre 5) et modèles de la méthode. Le « catalogue des
%% sujets du projet individuel » du squelette relève du cours INF3226 ; il
%% est remplacé ici par une grille d'auto-contrôle de l'observation.
%% ============================================================================

\chapter{Modèles et Grilles}
\label{ann:templates}

Cette annexe rassemble les outils de la méthode et les modèles à
reproduire en mission. Chacun renvoie au chapitre où il est expliqué.

\section*{Les outils O0 à O9}

\begin{tableaucours}{cL{3.6cm}L{1.4cm}L{6.4cm}}{%
    \tableauHeader{Outil} &
    \tableauHeader{Nom} &
    \tableauHeader{Phase} &
    \tableauHeader{Contenu}
}
O0 & Aide-mémoire du cycle & Toutes & Entrées, activités et livrables de
chaque phase \\
O1 & Fiche technique & 0 & Objectifs, acteurs, référentiels \\
O2 & Termes de référence & 0 & Quand, comment, par qui, avec quoi \\
O3 & Documents sollicités & 0 & Document, strate, priorité, interlocuteur,
date, cote \OUI{}/\NON{}/\NEK{} \\
O4 & Classeur d'aide à l'audit & 1 à 3 & Décomposition, fiches d'assertion,
cotation, $\mu$, $U$, $H$ \\
O5 & Guide d'entretien & 2 & Questions issues des fiches, par
interlocuteur \\
O6 & Plan d'audit & 1 & Contexte, périmètre, référentiels, risques, zones
retenues et écartées, programme, calendrier \\
O7 & Plan de tests & 2 & Code du test (DOC, ENT, OBS, TEST, MIX), preuve
attendue, papier de travail \\
O8 & Projets d'observation & 3 & Observations en sept éléments, soumises
au contradictoire \\
O9 & Suivi des documents & 2 & Mise à jour de O3 sur le terrain \\
\end{tableaucours}

\section*{La fiche d'assertion}

\begin{boxobjet}{Code · Intitulé du contrôle}
\textbf{Position~:} domaine $\rightarrow$ sous-domaine $\rightarrow$
processus $\rightarrow$ activité $\rightarrow$ contrôle.

\textbf{Objectif~:} « S'assurer que\ldots~» (à l'infinitif, sans « et » ni
« ou »).

\textbf{Critère~:} condition vérifiable, et sa source (procédure, loi,
référentiel).

\textbf{Question~:} celle que l'on cote \OUI{}, \NON{} ou \NEK{}, atomique.

\textbf{Risque~:} l'état de non-conformité si la réponse est \NON{}.

\textbf{Conséquence~:} ce qui arrive si le risque se réalise.

\textbf{Test, preuve, trace~:} code du test, pièce attendue, papier de
travail.
\end{boxobjet}

\section*{La règle de cotation}

\begin{tableaucours}{L{6.6cm}c}{%
    \tableauHeader{Situation} &
    \tableauHeader{Cote}
}
Le critère est satisfait, pièce à l'appui & \OUI{} \\
Le critère n'est pas satisfait, ou la pièce exigée est absente, ou le
document est périmé ou non approuvé & \NON{} \\
L'audité confirme par écrit que le document exigé n'existe pas & \NON{} \\
Le document n'est pas remis et l'on ne sait pas s'il existe & \NEK{} \\
Seule une déclaration orale est disponible & \NEK{} \\
Le contrôle est sans objet dans le périmètre & écarté, raison écrite \\
\end{tableaucours}

\section*{L'observation en sept éléments}

\begin{boxobjet}{Code · Observation}
\textbf{1. Norme~:} le critère et sa source.

\textbf{2. Constatation~:} le fait, chiffré, sa preuve, sa cote.

\textbf{3. Risques~:} ce qui peut arriver.

\textbf{4. Contradictoire~:} la réponse de l'audité, datée.

\textbf{5. Analyse} (niveau de confiance)~: la cause.

\textbf{6. Conclusion~:} la position finale de l'auditeur.

\textbf{7. Recommandations~:} ce qu'il faut faire, en traitant la cause~;
responsable, échéance, preuve d'exécution.
\end{boxobjet}

\section*{Grille d'auto-contrôle avant de rendre une observation}

\begin{tableaucours}{L{9.6cm}c}{%
    \tableauHeader{Question} &
    \tableauHeader{Oui~/ non}
}
Le critère est-il cité, avec sa source, et l'ai-je lu~? & \\
La constatation est-elle séparée de l'analyse~? & \\
Chaque chiffre est-il rattaché à un papier de travail et à une empreinte~? & \\
Les \NEK{} sont-ils signalés, avec ce qu'il faudrait pour les lever~? & \\
La réponse de l'audité figure-t-elle, datée~? & \\
Le texte évite-t-il toute qualification pénale~? & \\
La recommandation traite-t-elle la cause, avec un responsable et une
échéance~? & \\
Ai-je audité mon propre travail~? Si oui, l'ai-je déclaré~? & \\
\end{tableaucours}

\section*{Autres modèles}

\begin{itemize}
    \item \textbf{Lettre d'extraction}~: source, date et heure, requête,
          paramètres, volumes, empreinte, signataires
          (Chapitre~\ref{chap:preuve}).
    \item \textbf{Note d'autorisation d'un test technique}~: signataire
          habilité, périmètre, méthodes, fenêtres, contacts, sort des
          résultats (Chapitre~\ref{chap:intrusion}).
    \item \textbf{Grille de passage d'une mise en production}~: critères,
          mesures, organe, écarts acceptés, retour arrière
          (Chapitre~\ref{chap:audit-projets}).
    \item \textbf{Fiche d'indicateur}~: définition, formule, source,
          fréquence, responsable, cible (Chapitre~\ref{chap:audit-gouvernance}).
    \item \textbf{Registre des risques}~: identifiant, scénario, cause,
          vraisemblance, impact, score, contrôles, risque résiduel,
          propriétaire, traitement (Chapitre~\ref{chap:risque}).
    \item \textbf{Déclaration d'applicabilité}~: mesure, retenue ou exclue,
          justification, état de mise en œuvre, preuve
          (Chapitre~\ref{chap:audit-smsi}).
    \item \textbf{Registre de conservation d'une pièce}~: identifiant,
          source, acquisition, empreinte initiale, lieux, accès, remises
          (Chapitre~\ref{chap:audit-preuve}).
    \item \textbf{Procédure de remise à l'autorité judiciaire}~: geler,
          prouver l'intégrité, documenter la provenance, tenir le registre,
          remettre contre récépissé (Chapitre~\ref{chap:audit-preuve}).
    \item \textbf{Registre de conformité}~: obligation, source, contrôle,
          preuve, responsable (Chapitre~\ref{chap:audit-conformite}).
\end{itemize}
